\documentclass[12pt,letterpaper]{article}
\usepackage[utf8]{inputenc}
\usepackage[spanish,es-noshorthands]{babel}
\usepackage{amsmath,amssymb}
\usepackage{xcolor}
\usepackage{array,booktabs}
\usepackage{tikz}
\usetikzlibrary{calc,arrows.meta}
\usepackage[american]{circuitikz}
\usepackage[margin=2.2cm]{geometry}
\setlength{\parindent}{0pt}

% puntaje: numero en circulo rojo
\newcommand{\pts}[1]{\tikz[baseline=(c.base)]\node[circle,draw=red,text=red,thick,inner sep=1.3pt,font=\small\bfseries](c){#1};}
\newcommand{\uF}{\,\mu\mathrm{F}}
\newcommand{\uC}{\,\mu\mathrm{C}}
\newcommand{\Vv}{\,\mathrm{V}}

\begin{document}

\begin{center}
    {\Large\textbf{Pauta -- Certamen 3 Recuperativo -- Forma C}}\\[4pt]
    {\normalsize F\'isica II - Electromagnetismo \quad\textbullet\quad 09 de Noviembre, 2026}
\end{center}

\begin{center}
\renewcommand{\arraystretch}{1.2}
\begin{tabular}{lcc}
\toprule
Pregunta & Tema & Puntaje\\
\midrule
1 & Circuito de capacitores & 20\\
2 & Circuito resistivo (Kirchhoff) & 20\\
3 & Fuerza magn\'etica, momento y torque & 19\\
\midrule
 & \textbf{Total} & \textbf{59}\\
\bottomrule
\end{tabular}
\end{center}

%======================================================================
\section*{Pregunta 1 \quad\textcolor{red}{(20 puntos)}}

Datos: $C_1=3{,}0\uF$, $C_2=4{,}0\uF$, $C_3=3{,}0\uF$, $C_4=4{,}0\uF$, $C_5=2{,}0\uF$ y
$V_5=4{,}0\Vv$ (medido en $C_5$).

\subsection*{(a) Reducci\'on paso a paso \quad\textcolor{red}{(8 puntos)}}

\noindent\begin{minipage}[t]{0.58\linewidth}
\vspace{0pt}\textbf{Paso 1:} $C_4$ y $C_5$ comparten los mismos dos nodos: est\'an en \textbf{paralelo}.
\[
    C_{45} = C_4 + C_5 = 4{,}0\uF + 2{,}0\uF
    \;\Rightarrow\; \boxed{C_{45} = 6{,}0\uF}\ \pts{1}
\]
\end{minipage}\hfill
\begin{minipage}[t]{0.38\linewidth}
\vspace{0pt}\raggedleft\small Paso 1: \pts{1}\par
\begin{circuitikz}[scale=0.6]
    \draw (0,0) to[battery,invert,l=$V$] (0,3) to[C,l=$C_1$] (2.5,3);
    \draw (2.5,3) to[C,l=$C_2$] (2.5,0) to[short] (0,0);
    \draw (2.5,3) to[C,l=$C_3$] (5,3) to[C,l=$C_{45}$] (5,0) to[short] (2.5,0);
\end{circuitikz}
\end{minipage}\par\medskip

\noindent\begin{minipage}[t]{0.58\linewidth}
\vspace{0pt}\textbf{Paso 2:} $C_3$ y $C_{45}$ est\'an en \textbf{serie} (una sola rama):
\[
    C_{345} = \frac{C_3\,C_{45}}{C_3+C_{45}} = \frac{(3)(6)}{3+6}\uF
    \;\Rightarrow\; \boxed{C_{345} = 2{,}0\uF}\ \pts{1}
\]
\end{minipage}\hfill
\begin{minipage}[t]{0.38\linewidth}
\vspace{0pt}\raggedleft\small Paso 2: \pts{1}\par
\begin{circuitikz}[scale=0.6]
    \draw (0,0) to[battery,invert,l=$V$] (0,3) to[C,l=$C_1$] (2.5,3);
    \draw (2.5,3) to[C,l=$C_2$] (2.5,0) to[short] (0,0);
    \draw (2.5,3) to[short] (5,3) to[C,l=$C_{345}$] (5,0) to[short] (2.5,0);
\end{circuitikz}
\end{minipage}\par\medskip

\noindent\begin{minipage}[t]{0.58\linewidth}
\vspace{0pt}\textbf{Paso 3:} $C_2$ y $C_{345}$ est\'an en \textbf{paralelo}:
\begin{align*}
    C_{2345} &= C_2 + C_{345} = 4{,}0\uF + 2{,}0\uF\\
    \Rightarrow\ & \boxed{C_{2345} = 6{,}0\uF}\ \pts{1}
\end{align*}
\end{minipage}\hfill
\begin{minipage}[t]{0.38\linewidth}
\vspace{0pt}\raggedleft\small Paso 3:\par
\begin{circuitikz}[scale=0.6]
    \draw (0,0) to[battery,invert,l=$V$] (0,3) to[C,l=$C_1$] (3,3)
          to[C,l=$C_{2345}$] (3,0) to[short] (0,0);
\end{circuitikz}
\end{minipage}\par\medskip

\noindent\begin{minipage}[t]{0.58\linewidth}
\vspace{0pt}\textbf{Paso 4:} $C_1$ y $C_{2345}$ est\'an en \textbf{serie} con la fuente:
\begin{align*}
    C_{eq} &= \frac{C_1\,C_{2345}}{C_1+C_{2345}} = \frac{(3)(6)}{3+6}\uF\\
    \Rightarrow\ & \boxed{C_{eq} = 2{,}0\uF}\ \pts{2}
\end{align*}
\end{minipage}\hfill
\begin{minipage}[t]{0.38\linewidth}
\vspace{0pt}\raggedleft\small Capacitancia eq.\ final: \pts{1}\par
\begin{circuitikz}[scale=0.6]
    \draw (0,0) to[battery,invert,l=$V$] (0,3) to[short] (3,3)
          to[C,l=$C_{eq}$] (3,0) to[short] (0,0);
\end{circuitikz}
\end{minipage}\par

\subsection*{(b) Carga de cada capacitor \quad\textcolor{red}{(8 puntos)}}
\begin{itemize}\setlength\itemsep{4pt}
    \item $C_4$ y $C_5$ est\'an en paralelo, luego tienen el mismo voltaje $V_4 = V_5 = 4{,}0\Vv$:
    \[
        Q_5 = C_5 V_5 = (2{,}0\uF)(4{,}0\Vv) = 8{,}0\uC\ \pts{1}
        \qquad
        Q_4 = C_4 V_4 = (4{,}0\uF)(4{,}0\Vv) = 16\uC\ \pts{1}
    \]
    \item $C_3$ est\'a en serie con el bloque $C_{45}$, por lo que lleva su carga total:
    \[
        \boxed{Q_3 = Q_{45} = Q_4 + Q_5 = 16\uC + 8{,}0\uC = 24\uC}\ \pts{2}
    \]
    \item Voltaje en el bloque $C_{2345}$ ($C_2$ est\'a en paralelo con la rama $C_3$--$C_{45}$):
    \[
        V_3 = \frac{Q_3}{C_3} = \frac{24\uC}{3{,}0\uF} = 8{,}0\Vv,
        \qquad
        V_2 = V_3 + V_{45} = 8{,}0\Vv + 4{,}0\Vv = 12\Vv\ \pts{1}
    \]
    \item Carga de $C_2$:
    \[
        \boxed{Q_2 = C_2 V_2 = (4{,}0\uF)(12\Vv) = 48\uC}\ \pts{1}
    \]
    \item $C_1$ est\'a en serie con el bloque $C_{2345}$, por lo que lleva la carga total:
    \[
        \boxed{Q_1 = Q_2 + Q_3 = 48\uC + 24\uC = 72\uC}\ \pts{2}
    \]
\end{itemize}

\subsection*{(c) Voltaje de la fuente \quad\textcolor{red}{(4 puntos)}}
\[
    V_1 = \frac{Q_1}{C_1} = \frac{72\uC}{3{,}0\uF} = 24\Vv\ \pts{2}
    \qquad\Rightarrow\qquad
    \boxed{V = V_1 + V_2 = 24\Vv + 12\Vv = 36\Vv}\ \pts{2}
\]
\textit{Verificaci\'on:} $Q_1 = C_{eq}V = (2{,}0\uF)(36\Vv) = 72\uC$.

\newpage
%======================================================================
\section*{Pregunta 2 \quad\textcolor{red}{(20 puntos)}}

Datos: $R_1=R_4=4\,\Omega$, $R_2=6\,\Omega$, $R_3=2\,\Omega$, $R_5=12\,\Omega$, $V_1=24\Vv$,
$V_2=8\Vv$.

\begin{center}
\begin{circuitikz}[american,scale=.8,transform shape]
\draw (0,0) to[battery1,invert,l_=$V_1$] (0,3)
 to[R,l=$R_1$,i>^=$I_1$] (3,3) node[circ]{};
\draw (3,3) to[R,l=$R_2$,i>^=$I_2$] (3,0) node[circ]{} -- (0,0);
\draw (3,3) -- (4,3) node[circ]{} to[R,l=$R_3$] (6.5,3) to[R,l=$R_4$] (9,3) node[circ]{}
 -- (10,3) to[short,i>^=$I_3$] (10,0);
\draw (4,3) -- (4,1.5) to[R,l=$R_5$] (9,1.5) -- (9,3);
\draw (3,0) to[battery1,l=$V_2$] (10,0);
\node at (-.35,2.1) {$+$}; \node at (-.35,.9) {$-$};
\node at (5.9,-.35) {$+$}; \node at (7.1,-.35) {$-$};
\end{circuitikz}
\end{center}

\subsection*{(a) Reducci\'on de resistencias \quad\textcolor{red}{(5 puntos)}}
$R_3$ y $R_4$ est\'an en \textbf{serie} (misma corriente, en una sola rama):
\[
    \boxed{R_{34} = R_3 + R_4 = 2 + 4 = 6\,\Omega}\ \pts{2}
\]
$R_{34}$ y $R_5$ est\'an en \textbf{paralelo} (comparten los mismos dos nodos):
\[
    \frac{1}{R_{\mathrm{eq}}} = \frac{1}{R_{34}} + \frac{1}{R_5} = \frac{1}{6} + \frac{1}{12} = \frac{1}{4}
    \quad\Rightarrow\quad \boxed{R_{\mathrm{eq}} = 4\,\Omega}\ \pts{3}
\]

\subsection*{(b) Leyes de Kirchhoff \quad\textcolor{red}{(6 puntos)}}
Ley de nodos:
\begin{equation} I_1 = I_2 + I_3. \tag{1}\end{equation}
Malla 1 (izquierda):
\begin{equation} V_1 - I_1R_1 - I_2R_2 = 0
\quad\Rightarrow\quad 24 - 4I_1 - 6I_2 = 0. \tag{2}\end{equation}
Malla 2 (derecha, con $R_{\mathrm{eq}}$):
\begin{equation} V_2 + I_2R_2 - I_3R_{\mathrm{eq}} = 0
\quad\Rightarrow\quad 8 + 6I_2 - 4I_3 = 0. \tag{3}\end{equation}
De (3), $I_3 = 2 + 1{,}5\,I_2$; con (1), $I_1 = 2 + 2{,}5\,I_2$; reemplazando en (2):
$24 - 8 - 10I_2 - 6I_2 = 0 \Rightarrow I_2 = 1\ \mathrm{A}$. La soluci\'on del sistema es:
\[
    \boxed{I_1 = 4{,}5\ \mathrm{A}\ \pts{2},\qquad I_2 = 1{,}0\ \mathrm{A}\ \pts{2},\qquad
    I_3 = 3{,}5\ \mathrm{A}\ \pts{2}}
\]
Las tres corrientes resultan positivas en los sentidos se\~nalados.

\newpage
\subsection*{(c) Corriente, voltaje y potencia disipada \quad\textcolor{red}{(9 puntos)}}
Usando $\Delta V = IR$ y $P = I^2R$:
\begin{center}
\renewcommand{\arraystretch}{1.3}
\begin{tabular}{lccc}
\toprule
Elemento & $I\,(\mathrm{A})$ & $\Delta V\,(\mathrm{V})$ & $P\,(\mathrm{W})$\\
\midrule
$R_1 = 4\,\Omega$ & $4{,}5$ & $18$ & $81$\\
$R_2 = 6\,\Omega$ & $1{,}0$ & $6$ & $6$\\
$R_{\mathrm{eq}} = 4\,\Omega$ & $3{,}5$ & $14$ & $49$\\
\bottomrule
\end{tabular}

\smallskip
\textcolor{red}{(1 punto cada casilla)}
\end{center}
\textit{Verificaci\'on (malla 1):} $V_1 = \Delta V_{R_1} + \Delta V_{R_2} = 18 + 6 = 24\Vv$.

\newpage
%======================================================================
\section*{Pregunta 3 \quad\textcolor{red}{(19 puntos)}}

\begin{center}
\fbox{\begin{minipage}{0.9\textwidth}\centering
\textbf{Productos cruz de los vectores unitarios}\\[4pt]
$\hat{\imath}\times\hat{\jmath} = \hat{k},\quad \hat{\jmath}\times\hat{k} = \hat{\imath},\quad \hat{k}\times\hat{\imath} = \hat{\jmath}$
\qquad
$\hat{\jmath}\times\hat{\imath} = -\hat{k},\quad \hat{k}\times\hat{\jmath} = -\hat{\imath},\quad \hat{\imath}\times\hat{k} = -\hat{\jmath}$\\[3pt]
$\hat{\imath}\times\hat{\imath} = \hat{\jmath}\times\hat{\jmath} = \hat{k}\times\hat{k} = \vec{0}$
\end{minipage}}
\end{center}

\subsection*{3a) Carga puntual \quad\textcolor{red}{(7 puntos)}}

Datos: $q = +6\times10^{-6}\ [\text{C}]$, $v = 5\times10^{5}\ [\text{m/s}]$, $B = 0{,}04\ [\text{T}]$.
De la Fig.~3a: $\vec{v}$ forma $50^\circ$ sobre el eje $+x$ y $\vec{B}$ apunta $20^\circ$ sobre el eje $-x$:
\begin{align*}
    \vec{v} &= v\left(\cos50^\circ\,\hat{\imath} + \sin50^\circ\,\hat{\jmath}\right) && \pts{2}\\
    \vec{B} &= B\left(-\cos20^\circ\,\hat{\imath} + \sin20^\circ\,\hat{\jmath}\right) && \pts{2}
\end{align*}
Producto cruz t\'ermino a t\'ermino:
\begin{align*}
    \vec{v}\times\vec{B}
      &= vB\bigl[-\cos50^\circ\cos20^\circ(\hat{\imath}\times\hat{\imath})
         + \cos50^\circ\sin20^\circ(\hat{\imath}\times\hat{\jmath})\\
      &\qquad\;\; - \sin50^\circ\cos20^\circ(\hat{\jmath}\times\hat{\imath})
         + \sin50^\circ\sin20^\circ(\hat{\jmath}\times\hat{\jmath})\bigr]\\
      &= vB\left(\cos50^\circ\sin20^\circ + \sin50^\circ\cos20^\circ\right)\hat{k}
       = vB\sin70^\circ\,\hat{k}\\
      &= (5\times10^{5})(0{,}04)(0{,}9397)\,\hat{k} = 18794\,\hat{k}\ [\text{T m/s}]
\end{align*}
\[
    \vec{F}_m = q\,\vec{v}\times\vec{B} = (6\times10^{-6})(18794)\,\hat{k}
    \qquad\Rightarrow\qquad
    \boxed{\vec{F}_m = 1{,}13\times10^{-1}\,\hat{k}\ [\text{N}]} \quad\pts{3}
\]
La fuerza sale de la p\'agina (la carga es positiva). El \'angulo entre $\vec{v}$ y $\vec{B}$ es
$110^\circ$, y $\sin110^\circ = \sin70^\circ$.

\subsection*{3b) Lazo rectangular \quad\textcolor{red}{(12 puntos)}}

Datos: $a = 0{,}12\ [\text{m}]$, $b = 0{,}07\ [\text{m}]$, $I = 3{,}0\ [\text{A}]$ en sentido
\textbf{antihorario}, $\vec{B} = -0{,}4\,\hat{\imath} + 0{,}5\,\hat{k}\ [\text{T}]$.
Lazo de la Fig.~3b: $1=(0,0,0)$, $2=(a,0,0)$, $3=(a,b,0)$, $4=(0,b,0)$. Con el recorrido
$1\to2\to3\to4\to1$: $\vec{l}_{23} = b\,\hat{\jmath}$ (de $2$ a $3$) y $\vec{l}_{34} = -a\,\hat{\imath}$
(de $3$ a $4$); por la regla de la mano derecha $\hat{n} = +\hat{k}$.

\textbf{(i) Fuerzas} $\vec{F} = I\,\vec{l}\times\vec{B}$:
\begin{align*}
    \vec{F}_{23} &= (3{,}0)(0{,}07\,\hat{\jmath})\times(-0{,}4\,\hat{\imath} + 0{,}5\,\hat{k})
                  = 0{,}21\left[-0{,}4\,(\hat{\jmath}\times\hat{\imath}) + 0{,}5\,(\hat{\jmath}\times\hat{k})\right]\\
                 &= 0{,}21\left[0{,}4\,\hat{k} + 0{,}5\,\hat{\imath}\right]
\end{align*}
\[
    \boxed{\vec{F}_{23} = 0{,}105\,\hat{\imath} + 0{,}084\,\hat{k}\ [\text{N}]} \quad\pts{3}
\]
\begin{align*}
    \vec{F}_{34} &= (3{,}0)(-0{,}12\,\hat{\imath})\times(-0{,}4\,\hat{\imath} + 0{,}5\,\hat{k})
                  = -0{,}36\left[-0{,}4\,(\hat{\imath}\times\hat{\imath}) + 0{,}5\,(\hat{\imath}\times\hat{k})\right]\\
                 &= -0{,}36\left[0 - 0{,}5\,\hat{\jmath}\right]
\end{align*}
\[
    \boxed{\vec{F}_{34} = 0{,}18\,\hat{\jmath}\ [\text{N}]} \quad\pts{3}
\]

\textbf{(ii) Momento magn\'etico} (corriente antihoraria $\Rightarrow$ $\hat{n} = +\hat{k}$):
\[
    \vec{\mu} = I\,ab\,\hat{k} = (3{,}0)(0{,}12)(0{,}07)\,\hat{k}
    \qquad\Rightarrow\qquad
    \boxed{\vec{\mu} = 0{,}0252\,\hat{k}\ [\text{A m}^2]} \quad\pts{3}
\]

\textbf{(iii) Torque:}
\begin{align*}
    \vec{\tau} = \vec{\mu}\times\vec{B} &= 0{,}0252\,\hat{k}\times(-0{,}4\,\hat{\imath} + 0{,}5\,\hat{k})\\
               &= -0{,}01008\,(\hat{k}\times\hat{\imath}) + 0{,}0126\,(\hat{k}\times\hat{k})
\end{align*}
\[
    \boxed{\vec{\tau} = -1{,}01\times10^{-2}\,\hat{\jmath}\ [\text{N m}]} \quad\pts{3}
\]

\medskip
\noindent\textcolor{red}{\textbf{Puntaje total Pregunta 3: 19 puntos}} \textcolor{red}{(3a: 7 puntos; 3b: 12 puntos)}

\medskip
\textit{Nota:} sumando los cuatro tramos del lazo, la fuerza neta es $\vec{0}$, como corresponde a un
lazo cerrado en un campo uniforme ($\vec{F}_{12} = -\vec{F}_{34}$ y $\vec{F}_{41} = -\vec{F}_{23}$).

\end{document}
